%% ============================================================
%%  Section 2 — Literature Review (compact version)
%%  Article: BRKGA with Rolling Horizon for the MDRP
%%  Citation style: LNCS (numeric \cite{})
%% ============================================================

\section{Literature Review}\label{sec:litreview}

%% ------------------------------------------------------------------
\subsection{The Meal Delivery Routing Problem}\label{subsec:mdrp}
%% ------------------------------------------------------------------

The MDRP was formally defined by Reyes et al.~\cite{reyes2018mdrp},
who introduced a rolling-horizon framework in which orders arrive
stochastically, are picked up at restaurants, and are delivered by
couriers subject to time-window constraints.
Ulmer et al.~\cite{ulmer2021restaurant} extended this baseline by
modelling uncertain restaurant preparation times through the
\emph{Anticipative Customer Assignment} (ACA) model, which delays
courier--order matching to avoid premature assignments and long idle
waits—a design principle that informs the Business Model heuristic
presented in this paper.
Steever et al.~\cite{steever2019dynamic} addressed the closely
related Virtual Food Court Delivery Problem (VFCDP), using
continuous event-driven replanning, while Yildiz and
Savelsbergh~\cite{yildiz2019provably} proposed exact
column-generation methods for the deterministic variant, providing
high-quality benchmarks for evaluating heuristics.
More recent work has broadened the MDRP to include multi-objective
trade-offs~\cite{dusadeerungsikul2023multiobjective},
ghost kitchens~\cite{neria2024ghost} (delivery-only facilities with
no storefront that can be sited freely, adding spatial flexibility
to kitchen-and-route co-optimisation), and shared-logistics
modes~\cite{bi2023shared}; our work focuses on the single-platform,
heterogeneous-fleet setting dominant among Latin American operators.

%% ------------------------------------------------------------------
\subsection{Planning-Horizon Approaches}\label{subsec:horizon}
%% ------------------------------------------------------------------

MDRP solution strategies differ in their assumptions about future
information availability~\cite{reyes2018mdrp}.
\emph{Static methods} assume complete knowledge at the planning
horizon's start; they produce high-quality plans but cannot react
to new arrivals~\cite{dusadeerungsikul2023multiobjective,yildiz2019provably}.
\emph{Rolling-horizon} (RH) strategies re-solve a short-horizon
sub-problem at each time window, incorporating new
information~\cite{reyes2018mdrp,auad2023courier}.
Auad et al.~\cite{auad2023courier} dynamically resized courier
regions to improve coverage and satisfaction;
Mart\'{\i}nez-Sykora et al.~\cite{martinez2024fairness} applied RH
scheduling to balance fairness and efficiency;
Thipparthy et al.~\cite{thipparthy2024optimizing} showed that
batching decisions within each window substantially affect travel
distance and wait time.
\emph{Fully dynamic and stochastic} methods recalculate routes on
every event~\cite{steever2019dynamic} or incorporate probabilistic
demand forecasts~\cite{ulmer2021restaurant}; though most adaptive,
they rely on approximate solvers or learned
policies~\cite{bozanta2022reinforcement} due to high computational
demands.

%% ------------------------------------------------------------------
\subsection{Metaheuristics for Meal and Last-Mile Delivery}
\label{subsec:metaheuristics}
%% ------------------------------------------------------------------

Evolutionary algorithms have been applied successfully to MDRP
variants: Liao et al.~\cite{liao2020multiobjective} used NSGA-II to
balance workload and environmental impact in green meal delivery;
Xie and Wu~\cite{xie2024genetic} designed genetic operators to
enforce courier wage fairness without Lagrangian penalties;
Ba\~{n}os et al.~\cite{banos2013hybrid} hybridised a genetic
algorithm with local search for the multi-objective VRPTW,
establishing design principles—population diversity, elite
preservation, local refinement—that directly inform our BRKGA.
Reinforcement-learning approaches (Q-Learning,
DDQN~\cite{bozanta2022reinforcement}; deep
RL~\cite{ahamed2020deep}) show promise but require extensive
training data and offer limited interpretability.
Giraldo-Herrera and \'Alvarez-Mart\'{\i}nez~\cite{giraldo2024grasp}
proposed a GRASP for the MDRP on the same Latin American dataset
used here, providing a direct methodological baseline.
The BRKGA framework~\cite{goncalves2011brkga,londe2025review}
encodes solutions as random-key vectors in $[0,1]$ decoded into
feasible assignments without repair operators; biased crossover
($p_e=0.7$) concentrates search around elite solutions while
mutants maintain diversity. Its rapid per-window convergence makes
it well suited for rolling-horizon contexts, yet its integration
with an RH wrapper for a fully dynamic problem such as the MDRP
has not been previously studied.

%% ------------------------------------------------------------------
\subsection{Research Gap and Positioning}\label{subsec:gap}
%% ------------------------------------------------------------------

Two gaps motivate this work.
First, RH inner solvers are almost exclusively constructive greedy
heuristics; population-based search within each rolling window
has not been systematically evaluated.
Second, the quality--speed trade-off is rarely quantified across
heterogeneous instance sizes.
We address both by embedding a BRKGA inside an RH wrapper and
comparing two fitness variants (coverage-oriented and
travel-oriented) against a strong greedy baseline (Business Model)
on 24 real-world instances (538--59{,}167 orders).
Table~\ref{tab:litreview} summarises the positioning relative to
the closest related studies.

\begin{table}[ht]
\centering
\caption{Positioning of the present work in the literature.
RH: rolling horizon; Cov: coverage; TT: travel time.}
\label{tab:litreview}
\small
\setlength{\tabcolsep}{4pt}
\begin{tabular}{lcccccc}
\toprule
\textbf{Study} & \textbf{Method} & \textbf{RH} &
\textbf{Obj.} & \textbf{Dynamic} &
\textbf{Real data} & \textbf{Instance size} \\
\midrule
Reyes et al.~\cite{reyes2018mdrp}                        & Greedy       & \checkmark & TT       & \checkmark & \checkmark & Small--Med   \\
Ulmer et al.~\cite{ulmer2021restaurant}                  & DP/ACA       & \checkmark & TT       & \checkmark & \checkmark & Small--Med   \\
Steever et al.~\cite{steever2019dynamic}                 & Greedy       & \checkmark & TT       & \checkmark & \checkmark & Med          \\
Yildiz \& Savelsbergh~\cite{yildiz2019provably}          & Exact        &            & Cost     &            & \checkmark & Small--Med   \\
Liao et al.~\cite{liao2020multiobjective}                & NSGA-II      &            & Multi    &            & Synthetic  & Small        \\
Bozanta et al.~\cite{bozanta2022reinforcement}           & RL (DDQN)    & \checkmark & TT       & \checkmark & \checkmark & Small        \\
Auad et al.~\cite{auad2023courier}                       & Greedy RH    & \checkmark & Fairness & \checkmark & \checkmark & Med          \\
Giraldo-Herrera \& \'Alvarez~\cite{giraldo2024grasp}     & GRASP        &            & Cov      &            & \checkmark & Small--Large \\
Mart\'{\i}nez-Sykora et al.~\cite{martinez2024fairness}  & RH sched.    & \checkmark & Fairness & \checkmark & \checkmark & Small--Med   \\
\midrule
\textbf{This work}                                       & \textbf{BRKGA} & \checkmark & Cov/TT & \checkmark & \checkmark & Small--Large \\
\bottomrule
\end{tabular}
\end{table}
